\documentclass[12pt]{ximera}
\title{Exam 1 (2020)}
\begin{document}
\begin{abstract}
An archived exam on Chapter 1: majority thresholds, Plurality and Plurality-with-Elimination winners, fairness criteria, Condorcet candidates, a modified Borda count, and Pairwise Comparisons point totals.
\end{abstract}
\maketitle

\section*{Exam 1 (2020)}

\begin{problem}
Problems 1--3 refer to the following preference schedule for an election with four
candidates.
\begin{center}
\begin{tabular}{c|ccc}
Number of voters & 10 & 8 & 7\\\hline
1st & B & D & C\\
2nd & A & B & A\\
3rd & D & C & D\\
4th & C & A & B
\end{tabular}
\end{center}

\begin{prompt}
\textbf{1.} To obtain a majority of the first-place votes cast, a candidate needs

\begin{multipleChoice}
\choice{$25$ votes}
\choice{$12$ votes}
\choice[correct]{$13$ votes}
\choice{more votes than any other candidate}
\end{multipleChoice}

\begin{feedback}
There are $10+8+7=25$ voters, and more than half of $25$ is $13$.
\end{feedback}
\end{prompt}

\begin{prompt}
\textbf{2.} Which candidate wins under the Plurality method?

\begin{multipleChoice}
\choice{A}
\choice[correct]{B}
\choice{C}
\choice{D}
\choice{There is a tie}
\end{multipleChoice}

\begin{feedback}
First-place votes: B $10$, D $8$, C $7$, A $0$.
\end{feedback}
\end{prompt}

\begin{prompt}
\textbf{3.} Which candidate wins under Plurality-with-Elimination?

\begin{multipleChoice}
\choice{A}
\choice{B}
\choice{C}
\choice[correct]{D}
\choice{The method is not required / does not apply}
\end{multipleChoice}

\begin{feedback}
Eliminate A ($0$), then C ($7$). C's voters rank D next (A is gone), giving D $15$ and
B $10$. D wins.
\end{feedback}
\end{prompt}
\end{problem}

\begin{problem}
\begin{prompt}
\textbf{4.} True or false: the Condorcet criterion can be violated by the Pairwise
Comparisons method.

\begin{multipleChoice}
\choice{True}
\choice[correct]{False}
\end{multipleChoice}

\begin{feedback}
False. A Condorcet candidate wins every comparison, so always wins by Pairwise
Comparisons.
\end{feedback}
\end{prompt}

\begin{prompt}
\textbf{5.} True or false: the Majority criterion is always satisfied by the
Plurality-with-Elimination method.

\begin{multipleChoice}
\choice[correct]{True}
\choice{False}
\end{multipleChoice}

\begin{feedback}
True. A majority candidate wins in the first round, before anyone is eliminated.
\end{feedback}
\end{prompt}
\end{problem}

\begin{problem}
Problems 6--8 refer to the following preference schedule.
\begin{center}
\begin{tabular}{c|cccc}
Number of voters & 7 & 5 & 3 & 2\\\hline
1st & C & A & B & D\\
2nd & D & C & D & C\\
3rd & A & D & A & B\\
4th & B & B & C & A
\end{tabular}
\end{center}

\begin{prompt}
\textbf{6.} Use a modified Borda count with $6$ points for first place, $5$ for second,
$2$ for third, and $0$ for fourth to find a complete ranking.

A: $\answer{50}$ \quad B: $\answer{22}$ \quad C: $\answer{77}$ \quad D: $\answer{72}$

\smallskip
1st $\answer{C}$, 2nd $\answer{D}$, 3rd $\answer{A}$, 4th $\answer{B}$

\begin{feedback}
\[
\begin{aligned}
\text{A:}&\ 7\cdot2+5\cdot6+3\cdot2+2\cdot0=50\\
\text{B:}&\ 7\cdot0+5\cdot0+3\cdot6+2\cdot2=22\\
\text{C:}&\ 7\cdot6+5\cdot5+3\cdot0+2\cdot5=77\\
\text{D:}&\ 7\cdot5+5\cdot2+3\cdot5+2\cdot6=72
\end{aligned}
\]
(Check: $17$ ballots at $13$ points each is $221$.)
\end{feedback}
\end{prompt}

\begin{prompt}
\textbf{7.} Does this election have a Condorcet candidate?

\begin{multipleChoice}
\choice{Yes, A}
\choice{Yes, B}
\choice[correct]{Yes, C}
\choice{Yes, D}
\choice{No}
\end{multipleChoice}

\begin{feedback}
C beats A $9$--$8$, B $14$--$3$, and D $12$--$5$.
\end{feedback}
\end{prompt}

\begin{prompt}
\textbf{8.} Who wins under the Pairwise Comparisons method? $\answer{C}$

\begin{feedback}
The Condorcet candidate always wins by Pairwise Comparisons. The full tally is C $3$,
D $2$ (beating A $12$--$5$ and B $14$--$3$), A $1$ (beating B $12$--$5$), and B $0$.
\end{feedback}
\end{prompt}
\end{problem}

\begin{problem}
\textbf{9.} An election with six candidates (A--F) used the Pairwise Comparisons method.
Candidate A lost $3$ comparisons and tied $2$; B tied $2$; C lost $2$ and tied $2$; D lost
$2$; F lost $3$ and tied $1$.

\begin{prompt}
\textbf{(a)} How many points does each candidate get?

A: $\answer{1}$ \quad B: $\answer{4}$ \quad C: $\answer{2}$ \quad D: $\answer{3}$ \quad
E: $\answer{3.5}$ \quad F: $\answer{1.5}$

\begin{feedback}
Each candidate is in $5$ comparisons; a win is worth $1$ and a tie $\frac12$. So A has
$0$ wins and $2$ ties, B $3$ wins and $2$ ties, C $1$ win and $2$ ties, D $3$ wins, and
F $1$ win and $1$ tie. There are $\frac{6\cdot5}2=15$ points in all, so E has
$15-(1+4+2+3+1.5)=3.5$.
\end{feedback}
\end{prompt}

\begin{prompt}
\textbf{(b)} Give a complete ranking.

1st $\answer{B}$, 2nd $\answer{E}$, 3rd $\answer{D}$, 4th $\answer{C}$, 5th $\answer{F}$,
6th $\answer{A}$
\end{prompt}
\end{problem}

\end{document}
